\section{Measurement, photon and alignment controls}
\label{sec:s:budgets}

\paragraph{Dose convention.}
Each encoder's excitation is divided by its mean over the specimen analysis grid and the full 256-direction scan. This fixes the deposited excitation per complete scan. For $k$ directions, exposure is multiplied by $256/k$. Random subsets therefore match nominal dose in expectation, not exactly for each subset. The per-specimen exposure is set from the random plate's full-scan mean signal and reused for every encoder. Equal deposited dose does not force identical detected counts: illumination can correlate differently with the specimen. For the two learned plates, mean detected signal changes by about $-1$ to $-2$\,\% across the 12 crops, with individual changes approximately $-12$ to $+9$\,\%.

\paragraph{Measurement compression.}
\Cref{tab:s:compression} reports both depth metrics at $C=1,4,8$. These values count all camera samples, not camera images alone. At matched nominal dose, fewer directions mean longer exposure per image. The learned 32-image acquisition gives lower W1 than the 256-image random acquisition for TV and pretrained PnP, but the close neuron-trained-PnP means do not establish equivalence. The TV curve is not monotone in compression: validation can prefer a stronger regularization weight at $C=4$, which smooths neuronal features more than at $C=8$.

\begin{table*}[t]
\centering\small
\widetab{\begin{tabular}{@{}llcccccc@{}}
\toprule
Reconstruction & Plate & W1, $C=1$ & Centroid & W1, $C=4$ & Centroid & W1, $C=8$ & Centroid\\
\midrule
TV & Random & 0.249 & 0.187 & 0.298 & 0.223 & 0.318 & 0.243\\
TV & Learned, seed 0 & 0.170 & 0.135 & 0.208 & 0.166 & 0.196 & 0.153\\
TV & Learned, seed 1 & 0.160 & 0.124 & 0.202 & 0.158 & 0.195 & 0.151\\
Pretrained PnP & Random & 0.413 & 0.357 & 0.445 & 0.384 & 0.496 & 0.431\\
Pretrained PnP & Learned, seed 0 & 0.222 & 0.184 & 0.254 & 0.212 & 0.277 & 0.231\\
Pretrained PnP & Learned, seed 1 & 0.213 & 0.173 & 0.249 & 0.206 & 0.279 & 0.234\\
Neuron-trained PnP & Random & 0.203 & 0.161 & 0.234 & 0.184 & 0.282 & 0.221\\
Neuron-trained PnP & Learned, seed 0 & 0.142 & 0.110 & 0.167 & 0.130 & 0.206 & 0.163\\
\bottomrule
\end{tabular}
}
\caption{Measurement budgets on the 12 prespecified crops. W1 and centroid errors are in $\mu$m and come from the same reconstruction rows. $C=1$ uses all 256 directions; $C=4,8$ average five subsets of 64 and 32 directions. No missing independent-initialization results are imputed.}
\label{tab:s:compression}
\end{table*}

\paragraph{Noise conditions.}
The nominal model uses Poisson shot noise. The background condition adds a uniform mean equal to the random plate's mean expected signal for the same exposure, plus independent Gaussian read noise with standard deviation 2 electrons. The known mean background is subtracted before reconstruction. The same background level, exposure and noise protocol are used for each encoder. A separate low-dose condition reduces exposure to one tenth. Scalar reconstruction parameters are revalidated for each changed noise condition, while networks stay fixed. \Cref{tab:s:noise} reports all three methods.

\begin{table*}[t]
\centering\small
\widetab{\begin{tabular}{@{}llcccccc@{}}
\toprule
Condition & Reconstruction & W1, random & W1, learned & Centroid, random & Centroid, learned & W1 change & Centroid change\\
\midrule
Shot noise & TV & 0.318 & 0.196 & 0.243 & 0.153 & $-38$\,\% [$-43$, $-34$] & $-37$\,\% [$-41$, $-34$]\\
Shot noise & Pretrained PnP & 0.496 & 0.277 & 0.431 & 0.231 & $-44$\,\% [$-50$, $-39$] & $-46$\,\% [$-53$, $-41$]\\
Shot noise & Neuron-trained PnP & 0.282 & 0.206 & 0.221 & 0.163 & $-27$\,\% [$-34$, $-20$] & $-26$\,\% [$-35$, $-18$]\\
Background + read noise & TV & 0.324 & 0.227 & 0.249 & 0.180 & $-30$\,\% [$-38$, $-22$] & $-28$\,\% [$-36$, $-20$]\\
Background + read noise & Pretrained PnP & 0.627 & 0.297 & 0.548 & 0.247 & $-53$\,\% [$-57$, $-48$] & $-55$\,\% [$-60$, $-50$]\\
Background + read noise & Neuron-trained PnP & 0.288 & 0.210 & 0.228 & 0.166 & $-27$\,\% [$-35$, $-19$] & $-27$\,\% [$-36$, $-18$]\\
0.1 dose & TV & 0.466 & 0.292 & 0.365 & 0.236 & $-37$\,\% [$-45$, $-29$] & $-35$\,\% [$-44$, $-28$]\\
0.1 dose & Pretrained PnP & 0.834 & 0.407 & 0.733 & 0.352 & $-51$\,\% [$-57$, $-45$] & $-52$\,\% [$-58$, $-45$]\\
0.1 dose & Neuron-trained PnP & 0.345 & 0.239 & 0.269 & 0.186 & $-31$\,\% [$-37$, $-24$] & $-31$\,\% [$-38$, $-24$]\\
\bottomrule
\end{tabular}
}
\caption{Noise controls on 12 crops at $C=8$. Errors are in $\mu$m; changes and cell-bootstrap 95\,\% intervals compare the learned and random plates within the same condition.}
\label{tab:s:noise}
\end{table*}

\paragraph{Photon budget.}
The dose sweep uses factors $0.1$, $0.3$, $1$ and $3$ of the nominal exposure (\cref{tab:s:dose}). At dose $0.3$, learned-plate W1 is $0.234$, $0.302$ and $0.220\um$ for TV, pretrained PnP and neuron-trained PnP, compared with $0.318$, $0.496$ and $0.282\um$ for the random plate at dose 1. These reductions of 26\,\%, 39\,\% and 22\,\% describe the tested simulation and inference budgets. The random curve's flattening between doses 1 and 3 is not evidence of a fundamental photon limit.

\begin{table*}[t]
\centering\small
\begin{tabular}{@{}lccccc@{}}
\toprule
Reconstruction & Nominal dose & W1, random & W1, learned & Centroid, random & Centroid, learned\\
\midrule
TV & 0.1 & 0.466 & 0.292 & 0.365 & 0.236\\
TV & 0.3 & 0.442 & 0.234 & 0.341 & 0.185\\
TV & 1 & 0.318 & 0.196 & 0.243 & 0.153\\
TV & 3 & 0.315 & 0.192 & 0.242 & 0.151\\
Pretrained PnP & 0.1 & 0.834 & 0.407 & 0.733 & 0.352\\
Pretrained PnP & 0.3 & 0.631 & 0.302 & 0.550 & 0.251\\
Pretrained PnP & 1 & 0.496 & 0.277 & 0.431 & 0.230\\
Pretrained PnP & 3 & 0.494 & 0.276 & 0.430 & 0.230\\
Neuron-trained PnP & 0.1 & 0.345 & 0.239 & 0.269 & 0.186\\
Neuron-trained PnP & 0.3 & 0.304 & 0.220 & 0.236 & 0.174\\
Neuron-trained PnP & 1 & 0.282 & 0.207 & 0.221 & 0.164\\
Neuron-trained PnP & 3 & 0.275 & 0.203 & 0.218 & 0.162\\
\bottomrule
\end{tabular}

\caption{Dose sweep at $C=8$, 12 crops. W1 and centroid errors ($\mu$m) are computed from the same reconstructions. Each exposure factor is relative to the matched nominal-dose protocol.}
\label{tab:s:dose}
\end{table*}

\paragraph{Alignment.}
Distance perturbations set $D=6.1$ or $7.1\um$. In the characterized case, the operator is rebuilt from the known distance and scalar reconstruction parameters are revalidated; the phase is unchanged. In the nominal-operator case, reconstruction ignores the perturbation. Lateral error shifts the plate by one $0.32\um$ feature. Pointing errors are offsets of 0.005 or 0.01 in transverse direction cosine. All comparisons use neuron-trained PnP, $C=8$, the same 12 crops and five draws.

\Cref{tab:s:alignment} separates these assumptions. Known distance changes retain the optical gain. Ignoring distance or lateral registration can make the learned plate worse than the random one, including an 11\,\% W1 increase for a one-feature lateral shift. Rebuilding an operator with a known perturbation is not an experimentally demonstrated calibration procedure. Phase-quantization experiments are not used to claim a fabricable low-level plate: the simulated surround was not quantized with the trained center.

\begin{table*}[t]
\centering\small
\widetab{\begin{tabular}{@{}lcccccc@{}}
\toprule
Condition & W1, random & W1, learned & Centroid, random & Centroid, learned & W1 change & Centroid change\\
\midrule
none (design) & 0.282 & 0.206 & 0.221 & 0.163 & $-27$\,\% [$-34$, $-20$] & $-26$\,\% [$-35$, $-18$]\\
Known: D=6.1 & 0.285 & 0.205 & 0.225 & 0.162 & $-28$\,\% [$-35$, $-22$] & $-28$\,\% [$-37$, $-20$]\\
Known: D=7.1 & 0.284 & 0.209 & 0.226 & 0.166 & $-27$\,\% [$-34$, $-20$] & $-26$\,\% [$-35$, $-19$]\\
Nominal operator: scan tilt 0.005 & 0.318 & 0.254 & 0.253 & 0.196 & $-20$\,\% [$-25$, $-16$] & $-22$\,\% [$-30$, $-15$]\\
Nominal operator: scan tilt 0.01 & 0.379 & 0.328 & 0.308 & 0.251 & $-13$\,\% [$-16$, $-10$] & $-18$\,\% [$-22$, $-15$]\\
Nominal operator: D=6.1 & 0.563 & 0.586 & 0.528 & 0.571 & $+4$\,\% [$+0.3$, $+8$] & $+8$\,\% [$+3$, $+14$]\\
Nominal operator: D=7.1 & 0.447 & 0.473 & 0.383 & 0.438 & $+6$\,\% [$+5$, $+7$] & $+14$\,\% [$+12$, $+17$]\\
Nominal operator: lateral shift 1 feature & 0.451 & 0.500 & 0.372 & 0.400 & $+11$\,\% [$+2$, $+23$] & $+7$\,\% [$-0.4$, $+19$]\\
\bottomrule
\end{tabular}
}
\caption{Alignment controls with neuron-trained PnP. Errors are in $\mu$m; changes compare learned and random plates under the same perturbation. ``Known'' means the offset is supplied to operator assembly and validation; ``nominal operator'' means it is ignored. Positive changes favor the random plate.}
\label{tab:s:alignment}
\end{table*}
